\documentclass[11pt]{article}

% Change "review" to "final" to generate the final (sometimes called camera-ready) version.
% Change to "preprint" to generate a non-anonymous version with page numbers.
\usepackage[final]{acl}

% Standard package includes
\usepackage{times}
\usepackage{latexsym}

% For proper rendering and hyphenation of words containing Latin characters (including in bib files)
\usepackage[T1]{fontenc}
% For Vietnamese characters
% \usepackage[T5]{fontenc}
% See https://www.latex-project.org/help/documentation/encguide.pdf for other character sets

% This assumes your files are encoded as UTF8
\usepackage[utf8]{inputenc}

% This is not strictly necessary, and may be commented out,
% but it will improve the layout of the manuscript,
% and will typically save some space.
\usepackage{microtype}

% This is also not strictly necessary, and may be commented out.
% However, it will improve the aesthetics of text in
% the typewriter font.
\usepackage{inconsolata}

%Including images in your LaTeX document requires adding
%additional package(s)
\usepackage{graphicx}

% If the title and author information does not fit in the area allocated, uncomment the following
%
%\setlength\titlebox{<dim>}
%
% and set <dim> to something 5cm or larger.

\title{CAS 732: Replication Project Proposal}

% Author information can be set in various styles:
% For several authors from the same institution:
% \author{Author 1 \and ... \and Author n \\
%         Address line \\ ... \\ Address line}
% if the names do not fit well on one line use
%         Author 1 \\ {\bf Author 2} \\ ... \\ {\bf Author n} \\
% For authors from different institutions:
% \author{Author 1 \\ Address line \\  ... \\ Address line
%         \And  ... \And
%         Author n \\ Address line \\ ... \\ Address line}
% To start a separate ``row'' of authors use \AND, as in
% \author{Author 1 \\ Address line \\  ... \\ Address line
%         \AND
%         Author 2 \\ Address line \\ ... \\ Address line \And
%         Author 3 \\ Address line \\ ... \\ Address line}

\author{Samuel Khzym \\
  \texttt{khzyms@mcmaster.ca}}

%\author{
%  \textbf{First Author\textsuperscript{1}},
%  \textbf{Second Author\textsuperscript{1,2}},
%  \textbf{Third T. Author\textsuperscript{1}},
%  \textbf{Fourth Author\textsuperscript{1}},
% \\
%  \textbf{Fifth Author\textsuperscript{1,2}},
%  \textbf{Sixth Author\textsuperscript{1}},
%  \textbf{Seventh Author\textsuperscript{1}},
%  \textbf{Eighth Author \textsuperscript{1,2,3,4}},
% \\
%  \textbf{Ninth Author\textsuperscript{1}},
%  \textbf{Tenth Author\textsuperscript{1}},
%  \textbf{Eleventh E. Author\textsuperscript{1,2,3,4,5}},
%  \textbf{Twelfth Author\textsuperscript{1}},
% \\
%  \textbf{Thirteenth Author\textsuperscript{3}},
%  \textbf{Fourteenth F. Author\textsuperscript{2,4}},
%  \textbf{Fifteenth Author\textsuperscript{1}},
%  \textbf{Sixteenth Author\textsuperscript{1}},
% \\
%  \textbf{Seventeenth S. Author\textsuperscript{4,5}},
%  \textbf{Eighteenth Author\textsuperscript{3,4}},
%  \textbf{Nineteenth N. Author\textsuperscript{2,5}},
%  \textbf{Twentieth Author\textsuperscript{1}}
% \\
% \\
%  \textsuperscript{1}Affiliation 1,
%  \textsuperscript{2}Affiliation 2,
%  \textsuperscript{3}Affiliation 3,
%  \textsuperscript{4}Affiliation 4,
%  \textsuperscript{5}Affiliation 5
% \\
%  \small{
%    \textbf{Correspondence:} \href{mailto:email@domain}{email@domain}
%  }
% }

\begin{document}

\nolinenumbers

\maketitle

\section{Selected Paper}

The selected paper for replication is \textit{Transformer-XL: Attentive Language Models Beyond a Fixed-Length Context} [1], which augments the original transformer architecture with a recurrence mechanism. This allows it to overcome the limitation of losing information between segmented sections of the text, thus allowing for higher coherence over a longer context window. The code is publicly available at the following GitHub repo: \verb|github.com/kimiyoung/transformer-xl|.

The paper uses open datasets with longer contexts to train (WikiText-103, enwik8, text8, One Billion Word, and Penn Treebank), and demonstrates lower perplexity/bits-per-character scores than other state-of-the-art methods of the time. It also shows better relative effective context length (RECL), which is a paper-defined metric, compared to other methods.

\section{Datasets}

The paper used the following public datasets which will be used for reproduction. The paper used additional datasets, but the replication project will only endeavour to use these ones.

\begin{itemize}
\item WikiText-103: Set of over 100M tokens extracted from Wikipedia articles.
\item enwik8: Set of English Wikipedia articles.
\item text8: Set of English Wikipedia text.
\end{itemize}

\section{Models}

The paper trained several sizes of models, with the following names and number of parameters. Other sizes were trained, but the replication project will just attempt to train these ones.
\begin{itemize}
\item 12L Transformer-XL (41M). 
\item Transformer-XL Standard (151M).
\item 24L Transformer-XL (277M).
\end{itemize}

\section{Hypotheses}

The reproduction projection will attempt to reproduce the training in the paper on the same datasets and using the same metrics. The following hypotheses will be attempted to replicate, with scope potentially removed (ie. dropping larger model training) depending on ease of implementation and training.

\begin{itemize}
\item Acheive a bits-per-character (bpc) score of $\sim 1.06$ on the enwik8 dataset with the 41M parameter model (better than state-of-the-art of the time).
\item Acheive a perplexity score of $\sim 20.4$ on the WikiText-103 dataset with the 151M parameter model (better than state-of-the-art of the time).
\item Acheive a bits-per-character (bpc) score of $\sim 0.99$ on the enwik8 dataset with the 277M parameter model (better than state-of-the-art of the time).
\item Acheive a bits-per-character (bpc) score of $\sim 1.08$ on the text8 dataset with the 277M parameter model (better than state-of-the-art of the time).
\end{itemize}

\section{Predicted Challenges}

Since the public GitHub repo provides the code and scripts to download the same datasets used in the paper, I believe getting the same data used by the authors should not be a concern. The main concern I think will be a challenge will be the training phase, since I do not have access to much compute (NVIDIA GTX 3070Ti), and I'm afraid the larger models might not converge quickly enough to acheive the results presented in the paper.

\section{References}

[1] Z. Dai, Z. Yang, Y. Yang, J. Carbonell, Q. V. Le, and R. Salakhutdinov, “Transformer-XL: Attentive Language Models Beyond a Fixed-Length Context,” arXiv.org, 2019. https://arxiv.org/abs/1901.02860

\end{document}
